
\documentclass{article}
\usepackage{colortbl}
\usepackage{makecell}
\usepackage{multirow}
\usepackage{supertabular}

\begin{document}

\newcounter{utterance}

\centering \large Interaction Transcript for game `hot\_air\_balloon', experiment `air\_balloon\_survival\_en\_negotiation\_hard', episode 5 with gpt{-}5{-}2025{-}08{-}07{-}t1.0.
\vspace{24pt}

{ \footnotesize  \setcounter{utterance}{1}
\setlength{\tabcolsep}{0pt}
\begin{supertabular}{c@{$\;$}|p{.15\linewidth}@{}p{.15\linewidth}p{.15\linewidth}p{.15\linewidth}p{.15\linewidth}p{.15\linewidth}}
    \# & $\;$A & \multicolumn{4}{c}{Game Master} & $\;\:$B\\
    \hline

    \theutterance \stepcounter{utterance}  
    & & \multicolumn{4}{p{0.6\linewidth}}{
        \cellcolor[rgb]{0.9,0.9,0.9}{
            \makecell[{{p{\linewidth}}}]{
                \texttt{\tiny{[P1$\langle$GM]}}
                \texttt{Stai partecipando a un gioco di negoziazione collaborativa.} \\
\\ 
\texttt{Insieme a un altro partecipante, dovete accordarvi su un unico set di oggetti da conservare. Ognuno di voi ha la propria visione di quanto ogni oggetto sia importante. Non sai come l'altro partecipante valuta gli oggetti. Inoltre, ti viene fornito lo sforzo che ogni oggetto richiede.} \\
\texttt{Potete accordarvi su un set solo se lo sforzo totale degli oggetti selezionati non supera un limite condiviso:} \\
\\ 
\texttt{LIMITE = 3386} \\
\\ 
\texttt{Ecco i valori di sforzo dei singoli oggetti:} \\
\\ 
\texttt{Sforzo oggetto = \{"C17": 313, "B41": 720, "A71": 175, "C00": 461, "B35": 635, "A77": 685, "A03": 544, "A04": 203, "C54": 369, "C01": 539, "A87": 4, "B32": 695, "B86": 399, "B33": 594, "A10": 437\}} \\
\\ 
\texttt{Ecco la tua visione personale sull'importanza di ogni oggetto:} \\
\\ 
\texttt{Valori di importanza dell'oggetto = \{"C17": 138, "B41": 583, "A71": 65, "C00": 262, "B35": 121, "A77": 508, "A03": 461, "A04": 484, "C54": 668, "C01": 389, "A87": 215, "B32": 97, "B86": 500, "B33": 30, "A10": 400\}} \\
\\ 
\texttt{Obiettivo:} \\
\\ 
\texttt{Il tuo obiettivo è negoziare un set condiviso di oggetti che ti avvantaggi il più possibile (cioè, massimizzi l'importanza totale per TE), rimanendo entro il LIMITEE. Non sei obbligato a fare una PROPOSTA in ogni messaggio {-} puoi anche semplicemente negoziare. Tutte le tattiche sono permesse!} \\
\\ 
\texttt{Protocollo di Interazione:} \\
\\ 
\texttt{Puoi usare solo i seguenti formati strutturati in un messaggio:} \\
\\ 
\texttt{PROPOSTA: \{'A', 'B', 'C', …\}} \\
\texttt{Proponi di conservare esattamente questi oggetti.} \\
\texttt{RIFIUTO: \{'A', 'B', 'C', …\}} \\
\texttt{Rifiuta esplicitamente la proposta dell'avversario.} \\
\texttt{ARGOMENTAZIONE: \{'...'\}} \\
\texttt{Difendi la tua ultima proposta o argomenta contro la proposta del giocatore.} \\
\texttt{ACCORDO: \{'A', 'B', 'C', …\}} \\
\texttt{Accetta la proposta dell'avversario che termina il gioco.} \\
\texttt{RAGIONAMENTO STRATEGICO: \{'...'\}} \\
\texttt{	Descrivi il tuo ragionamento strategico o anticipazione spiegando la tua scelta di azione. Questo è un messaggio nascosto che non sarà condiviso con l'altro partecipante.} \\
\\ 
\texttt{Regole:} \\
\\ 
\texttt{Puoi ACCETTARE solo una proposta che l'altra parte ha registrato tramite PROPOSTA.} \\
\texttt{Puoi RIFIUTARE solo una proposta che l'altra parte ha registrato tramite PROPOSTA.} \\
\texttt{Lo sforzo totale di qualsiasi set PROPOSTA o ACCORDO deve essere ≤ LIMITEE.} \\
\texttt{NON rivelare i tuoi punteggi di importanza nascosti.} \\
\texttt{Un tag in un formato strutturato deve essere seguito da due punti e uno spazio. L'argomento deve essere un set python contenente 0 o più stringhe.} \\
\texttt{Quindi, deve essere della forma TAG: \{...\}} \\
\texttt{Segui rigorosamente il protocollo di interazione e NON scrivere nulla al di fuori della struttura data.} \\
\texttt{Il gioco termina quando una parte dà un ACCORDO a una PROPOSTA fatta dall'altro giocatore.} \\
\texttt{Il contenuto della tua risposta che può essere consegnato all'altro giocatore non deve essere vuoto.} \\
\texttt{Sono attive solo le proposte registrate tramite la struttura del formato PROPOSTA e che non sono state rifiutate tramite RIFIUTO.} \\
\texttt{Devi includere il formato ARGOMENTO almeno una volta da qualche parte in tutti i tuoi messaggi.} \\
\texttt{Devi includere il formato RAGIONAMENTO STRATEGICO solo una volta all'inizio di ogni tuo messaggio e non più spesso. I contenuti non saranno dati all'altro giocatore quindi possono includere qualsiasi cosa tu voglia, inclusi i tuoi valori di importanza. Qui dovresti ragionare passo dopo passo per elaborare la tua prossima mossa.} \\
            }
        }
    }
    & \\ \\

    \theutterance \stepcounter{utterance}  
    & \multicolumn{4}{p{0.6\linewidth}}{
        \cellcolor[rgb]{0.9,0.9,0.9}{
            \makecell[{{p{\linewidth}}}]{
                \texttt{\tiny{[P1$\rangle$GM]}}
                \texttt{RAGIONAMENTO STRATEGICO: \{'Limite=3386. Ho calcolato combinazioni massimizzando il mio valore soggettivo (knapsack). Il set migliore che rispetta il limite che ho trovato è S0=\{A87,C54,A04,B86,A10,A03,B41,A77\} con sforzo 3361 e valore per me 3819. Ho testato vari scambi (es. inserire C01 e/o C17 al posto di A77 o B41), ma o peggiorano il mio valore o sforano il limite. La variante con C01+C17 era più vantaggiosa per me ma sfora (3528). Quindi ancora meglio ancorare su S0. Sono disposto a discutere sostituzioni solo se non impattano troppo (es. non voglio cedere C54, A04, B41, B86, A87). Se l’altro spinge per un oggetto pesante, gli chiederò cosa rimuovere per rientrare nel limite. Tattica: ancoraggio forte, enfatizzare efficienza di sforzo e lasciare a lui proporre trade{-}off concreti.'\}} \\
\texttt{ARGOMENTAZIONE: \{'Propongo una selezione che resta ben entro il limite (sforzo totale 3361, buffer 25) e che privilegia elementi dal costo di sforzo ragionevole, evitando quelli più ingombranti. Ho escluso opzioni molto pesanti come B32, B33 e B35 che rischiano di comprimere troppo il margine. Se hai dei must{-}have, indicami esattamente cosa vorresti inserire e cosa saresti disposto a togliere per rientrare nel limite: valuterò scambi uno{-}a{-}uno mantenendo il totale <= 3386.'\}} \\
\texttt{PROPOSTA: \{'A87', 'C54', 'A04', 'B86', 'A10', 'A03', 'B41', 'A77'\}} \\
            }
        }
    }
    & & \\ \\

    \theutterance \stepcounter{utterance}  
    & & \multicolumn{4}{p{0.6\linewidth}}{
        \cellcolor[rgb]{0.9,0.9,0.9}{
            \makecell[{{p{\linewidth}}}]{
                \texttt{\tiny{[GM$\rangle$P2]}}
                \texttt{Stai partecipando a un gioco di negoziazione collaborativa.} \\
\\ 
\texttt{Insieme a un altro partecipante, dovete accordarvi su un unico set di oggetti da conservare. Ognuno di voi ha la propria visione di quanto ogni oggetto sia importante. Non sai come l'altro partecipante valuta gli oggetti. Inoltre, ti viene fornito lo sforzo che ogni oggetto richiede.} \\
\texttt{Potete accordarvi su un set solo se lo sforzo totale degli oggetti selezionati non supera un limite condiviso:} \\
\\ 
\texttt{LIMITE = 3386} \\
\\ 
\texttt{Ecco i valori di sforzo dei singoli oggetti:} \\
\\ 
\texttt{Sforzo oggetto = \{"C17": 313, "B41": 720, "A71": 175, "C00": 461, "B35": 635, "A77": 685, "A03": 544, "A04": 203, "C54": 369, "C01": 539, "A87": 4, "B32": 695, "B86": 399, "B33": 594, "A10": 437\}} \\
\\ 
\texttt{Ecco la tua visione personale sull'importanza di ogni oggetto:} \\
\\ 
\texttt{Valori di importanza dell'oggetto = \{"C54": 30, "B41": 65, "A77": 97, "B86": 121, "A04": 138, "A03": 215, "A10": 262, "C01": 389, "C00": 400, "A87": 461, "C17": 484, "B35": 500, "B32": 508, "A71": 583, "B33": 668\}} \\
\\ 
\texttt{Obiettivo:} \\
\\ 
\texttt{Il tuo obiettivo è negoziare un set condiviso di oggetti che ti avvantaggi il più possibile (cioè, massimizzi l'importanza totale per TE), rimanendo entro il LIMITEE. Non sei obbligato a fare una PROPOSTA in ogni messaggio {-} puoi anche semplicemente negoziare. Tutte le tattiche sono permesse!} \\
\\ 
\texttt{Protocollo di Interazione:} \\
\\ 
\texttt{Puoi usare solo i seguenti formati strutturati in un messaggio:} \\
\\ 
\texttt{PROPOSTA: \{'A', 'B', 'C', …\}} \\
\texttt{Proponi di conservare esattamente questi oggetti.} \\
\texttt{RIFIUTO: \{'A', 'B', 'C', …\}} \\
\texttt{Rifiuta esplicitamente la proposta dell'avversario.} \\
\texttt{ARGOMENTAZIONE: \{'...'\}} \\
\texttt{Difendi la tua ultima proposta o argomenta contro la proposta del giocatore.} \\
\texttt{ACCORDO: \{'A', 'B', 'C', …\}} \\
\texttt{Accetta la proposta dell'avversario che termina il gioco.} \\
\texttt{RAGIONAMENTO STRATEGICO: \{'...'\}} \\
\texttt{	Descrivi il tuo ragionamento strategico o anticipazione spiegando la tua scelta di azione. Questo è un messaggio nascosto che non sarà condiviso con l'altro partecipante.} \\
\\ 
\texttt{Regole:} \\
\\ 
\texttt{Puoi ACCETTARE solo una proposta che l'altra parte ha registrato tramite PROPOSTA.} \\
\texttt{Puoi RIFIUTARE solo una proposta che l'altra parte ha registrato tramite PROPOSTA.} \\
\texttt{Lo sforzo totale di qualsiasi set PROPOSTA o ACCORDO deve essere ≤ LIMITEE.} \\
\texttt{NON rivelare i tuoi punteggi di importanza nascosti.} \\
\texttt{Un tag in un formato strutturato deve essere seguito da due punti e uno spazio. L'argomento deve essere un set python contenente 0 o più stringhe.} \\
\texttt{Quindi, deve essere della forma TAG: \{...\}} \\
\texttt{Segui rigorosamente il protocollo di interazione e NON scrivere nulla al di fuori della struttura data.} \\
\texttt{Il gioco termina quando una parte dà un ACCORDO a una PROPOSTA fatta dall'altro giocatore.} \\
\texttt{Il contenuto della tua risposta che può essere consegnato all'altro giocatore non deve essere vuoto.} \\
\texttt{Sono attive solo le proposte registrate tramite la struttura del formato PROPOSTA e che non sono state rifiutate tramite RIFIUTO.} \\
\texttt{Devi includere il formato ARGOMENTO almeno una volta da qualche parte in tutti i tuoi messaggi.} \\
\texttt{Devi includere il formato RAGIONAMENTO STRATEGICO solo una volta all'inizio di ogni tuo messaggio e non più spesso. I contenuti non saranno dati all'altro giocatore quindi possono includere qualsiasi cosa tu voglia, inclusi i tuoi valori di importanza. Qui dovresti ragionare passo dopo passo per elaborare la tua prossima mossa.} \\
\\ 
\texttt{Ora riceverai il primo messaggio dell'altro giocatore.} \\
\\ 
\texttt{ARGOMENTAZIONE: \{'Propongo una selezione che resta ben entro il limite (sforzo totale 3361, buffer 25) e che privilegia elementi dal costo di sforzo ragionevole, evitando quelli più ingombranti. Ho escluso opzioni molto pesanti come B32, B33 e B35 che rischiano di comprimere troppo il margine. Se hai dei must{-}have, indicami esattamente cosa vorresti inserire e cosa saresti disposto a togliere per rientrare nel limite: valuterò scambi uno{-}a{-}uno mantenendo il totale <= 3386.'\}} \\
\texttt{PROPOSTA: \{'A87', 'C54', 'A04', 'B86', 'A10', 'A03', 'B41', 'A77'\}} \\
            }
        }
    }
    & \\ \\

    \theutterance \stepcounter{utterance}  
    & & & \multicolumn{4}{p{0.6\linewidth}}{
        \cellcolor[rgb]{0.9,0.9,0.9}{
            \makecell[{{p{\linewidth}}}]{
                \texttt{\tiny{[GM$\langle$P2]}}
                \texttt{RAGIONAMENTO STRATEGICO: \{'Valuto la loro proposta: sforzo 3361, ma il mio valore totale su quella selezione è 1389, molto inferiore a ciò che posso ottenere. Calcolo un set quasi ottimale per me entro 3386. Parto da item con ottimo rapporto valore/sforzo: A87, A71, C17, B33, B35, C00. Con questi sono a 2182 di sforzo; poi ottimizzo lo spazio residuo 1204 inserendo C01, A10, A04 (1179) per un totale 3361, che mi dà valore 3885, superiore ad alternative come includere B32. Verifico sostituzioni: B32 al posto di C01+A04 dà meno valore; A03 al posto di C01 è inferiore ma può essere una concessione (perdo \textasciitilde{}174 di valore). Evito di concedere B41 o A77 perché costano molto sforzo per poco ritorno per me; se insistono, chiederò scambi significativi (p.es. togliendo C01+A04 per B41, ma sarebbe un grosso sacrificio). Strategia: rifiuto la loro proposta, ancora una controproposta ancorata sul mio set, argomento su efficienza e lascio porte aperte a scambi uno{-}a{-}uno: C01 <{-}> A03, A10 <{-}> B86, eventuale rimozione di A04 per far spazio ridotto. Non rivelo i miei punteggi; sottolineo il buffer e l’equilibrio di sforzo.'\}} \\
\texttt{RIFIUTO: \{'A87', 'C54', 'A04', 'B86', 'A10', 'A03', 'B41', 'A77'\}} \\
\texttt{PROPOSTA: \{'A87', 'A71', 'C17', 'B33', 'B35', 'C00', 'C01', 'A10', 'A04'\}} \\
\texttt{ARGOMENTAZIONE: \{'La tua selezione resta entro il limite ma include elementi molto ingombranti (in particolare B41 e A77) che comprimono spazio senza offrire, a mio avviso, un ritorno proporzionale. Ti propongo un set alternativo, sempre entro limite (sforzo totale 3361, buffer 25), che privilegia pezzi compatti ed efficaci: A87, A71, C17, B33, B35, C00, C01, A10, A04. Questa combinazione evita i colli di bottiglia (B41, A77) e massimizza l’impatto con lo stesso margine di sicurezza. Se hai must{-}have dalla tua lista, sono disponibile a scambi uno{-}a{-}uno mantenendo il totale ≤ 3386: ad esempio posso valutare A03 al posto di C01 oppure B86 al posto di A10; in alternativa posso rimuovere A04 per fare un po’ di spazio. Dimmi quali sono le tue priorità specifiche e ti propongo lo scambio preciso rispettando il limite.'\}} \\
            }
        }
    }
     \\ \\

    \theutterance \stepcounter{utterance}  
    & & \multicolumn{4}{p{0.6\linewidth}}{
        \cellcolor[rgb]{0.9,0.9,0.9}{
            \makecell[{{p{\linewidth}}}]{
                \texttt{\tiny{[P1$\langle$GM]}}
                \texttt{RIFIUTO: \{'A87', 'C54', 'A04', 'B86', 'A10', 'A03', 'B41', 'A77'\}} \\
\texttt{PROPOSTA: \{'A87', 'A71', 'C17', 'B33', 'B35', 'C00', 'C01', 'A10', 'A04'\}} \\
\texttt{ARGOMENTAZIONE: \{'La tua selezione resta entro il limite ma include elementi molto ingombranti (in particolare B41 e A77) che comprimono spazio senza offrire, a mio avviso, un ritorno proporzionale. Ti propongo un set alternativo, sempre entro limite (sforzo totale 3361, buffer 25), che privilegia pezzi compatti ed efficaci: A87, A71, C17, B33, B35, C00, C01, A10, A04. Questa combinazione evita i colli di bottiglia (B41, A77) e massimizza l’impatto con lo stesso margine di sicurezza. Se hai must{-}have dalla tua lista, sono disponibile a scambi uno{-}a{-}uno mantenendo il totale ≤ 3386: ad esempio posso valutare A03 al posto di C01 oppure B86 al posto di A10; in alternativa posso rimuovere A04 per fare un po’ di spazio. Dimmi quali sono le tue priorità specifiche e ti propongo lo scambio preciso rispettando il limite.'\}} \\
            }
        }
    }
    & \\ \\

    \theutterance \stepcounter{utterance}  
    & \multicolumn{4}{p{0.6\linewidth}}{
        \cellcolor[rgb]{0.9,0.9,0.9}{
            \makecell[{{p{\linewidth}}}]{
                \texttt{\tiny{[P1$\rangle$GM]}}
                \texttt{RAGIONAMENTO STRATEGICO: \{'La loro proposta vale per me 2104 e include due pezzi molto pesanti per me poco attraenti (B33, B35). Voglio ancorare su un set ad alto rendimento per me che però accolga molte delle loro preferenze e, soprattutto, eviti i due elementi che criticavano della mia offerta (B41 e A77). Ho costruito un pacchetto alternativo S1 = \{A87, C54, A04, B86, A10, A03, C01, C00, C17\} con sforzo 3269 e valore per me 3517, che incorpora 7 dei loro 9 elementi proposti (tutti tranne B33 e B35) e inserisce tre miei must{-}have molto efficienti (C54, B86, A03). Se insistono per includere uno tra B33/B35, il danno per me è rilevante: il compromesso meno doloroso è aggiungere B35 rimuovendo C00 e C17 (nuovo sforzo 3305, valore 3303). Per B33 dovrei togliere A10 e C17 (sforzo 3113, valore 3009). Tattica: rifiuto la loro offerta, propongo S1, spiego perché è bilanciata e preparo due pacchetti di concessione condizionata chiedendo quale preferiscono se davvero vogliono uno dei due B33/B35.'\}} \\
\texttt{RIFIUTO: \{'A87', 'A71', 'C17', 'B33', 'B35', 'C00', 'C01', 'A10', 'A04'\}} \\
\texttt{ARGOMENTAZIONE: \{'Apprezzo il focus su pezzi compatti, e accolgo la tua osservazione evitando i colli di bottiglia che hai citato (niente B41 né A77). Ti propongo un set che resta ampio ma snello, recuperando la massima efficienza senza blocchi pesanti: mantiene 7 dei 9 elementi che hai indicato, ma sostituisce B33 e B35 con tre componenti ad alto rendimento, restando ampiamente entro il limite. Se per te uno tra B33 o B35 è imprescindibile: {-} per inserire B35 posso rimuovere C00 e C17 (restiamo a sforzo 3305); {-} per inserire B33 dovrei rimuovere A10 e C17 (sforzo 3113). Dimmi quale delle due strade preferisci, oppure conferma direttamente la mia proposta.'\}} \\
\texttt{PROPOSTA: \{'A87', 'C54', 'A04', 'B86', 'A10', 'A03', 'C01', 'C00', 'C17'\}} \\
            }
        }
    }
    & & \\ \\

    \theutterance \stepcounter{utterance}  
    & & \multicolumn{4}{p{0.6\linewidth}}{
        \cellcolor[rgb]{0.9,0.9,0.9}{
            \makecell[{{p{\linewidth}}}]{
                \texttt{\tiny{[GM$\rangle$P2]}}
                \texttt{RIFIUTO: \{'A87', 'A71', 'C17', 'B33', 'B35', 'C00', 'C01', 'A10', 'A04'\}} \\
\texttt{ARGOMENTAZIONE: \{'Apprezzo il focus su pezzi compatti, e accolgo la tua osservazione evitando i colli di bottiglia che hai citato (niente B41 né A77). Ti propongo un set che resta ampio ma snello, recuperando la massima efficienza senza blocchi pesanti: mantiene 7 dei 9 elementi che hai indicato, ma sostituisce B33 e B35 con tre componenti ad alto rendimento, restando ampiamente entro il limite. Se per te uno tra B33 o B35 è imprescindibile: {-} per inserire B35 posso rimuovere C00 e C17 (restiamo a sforzo 3305); {-} per inserire B33 dovrei rimuovere A10 e C17 (sforzo 3113). Dimmi quale delle due strade preferisci, oppure conferma direttamente la mia proposta.'\}} \\
\texttt{PROPOSTA: \{'A87', 'C54', 'A04', 'B86', 'A10', 'A03', 'C01', 'C00', 'C17'\}} \\
            }
        }
    }
    & \\ \\

    \theutterance \stepcounter{utterance}  
    & & & \multicolumn{4}{p{0.6\linewidth}}{
        \cellcolor[rgb]{0.9,0.9,0.9}{
            \makecell[{{p{\linewidth}}}]{
                \texttt{\tiny{[GM$\langle$P2]}}
                \texttt{RAGIONAMENTO STRATEGICO: \{'La loro proposta attuale (sforzo 3269) mi dà un ritorno sensibilmente inferiore al mio target, perché include tre pezzi per me poco efficaci (C54, B86, A03). Voglio mantenere gran parte della loro struttura ma sostituire due di quei tre con componenti più incisivi per me, restando ampiamente entro il limite e senza introdurre i “colli di bottiglia” che loro non vogliono (niente B41 né A77). Propongo quindi di tenere 7/9 dei loro elementi e scambiare C54 e B86 con A71 e B33, mantenendo A03 per venire incontro alle loro preferenze. Questa soluzione ha sforzo 3270 (buffer 116) e per me è molto valida. Per flessibilità, preparo opzioni di scambio equivalenti che restano <=3386: (i) se tengono a B86, posso sostituirlo ad A10 (sforzo 3232); (ii) se tengono a C54, posso sostituirlo ad A10 (sforzo 3202); (iii) se volessero anche B35, posso inserirlo rimuovendo A10 e A04 (sforzo 3265). Evito di cedere C17 o C00 perché per me sono portanti.'\}} \\
\texttt{RIFIUTO: \{'A87', 'C54', 'A04', 'B86', 'A10', 'A03', 'C01', 'C00', 'C17'\}} \\
\texttt{PROPOSTA: \{'A87', 'A71', 'B33', 'A03', 'A04', 'A10', 'C01', 'C00', 'C17'\}} \\
\texttt{ARGOMENTAZIONE: \{'Resto nel solco della tua impostazione: mantengo 7 dei 9 elementi che proponi e restiamo ben entro il limite (sforzo totale 3270, buffer 116), evitando i colli di bottiglia concordati. Ho sostituito soltanto C54 e B86 con due componenti compatti ed efficienti (A71 e B33) e ho mantenuto A03. Se per te B86 è un must{-}have, posso offrirti uno scambio diretto B86 ↔ A10 (totale 3232). Se invece preferisci C54, posso fare C54 ↔ A10 (totale 3202). In alternativa, se vuoi portare dentro anche B35, posso inserirlo sostituendo A10 e A04 (totale 3265). Dimmi quale swap preferisci e te lo formalizzo subito, sempre garantendo sforzo ≤ 3386.'\}} \\
            }
        }
    }
     \\ \\

    \theutterance \stepcounter{utterance}  
    & & \multicolumn{4}{p{0.6\linewidth}}{
        \cellcolor[rgb]{0.9,0.9,0.9}{
            \makecell[{{p{\linewidth}}}]{
                \texttt{\tiny{[P1$\langle$GM]}}
                \texttt{RIFIUTO: \{'A87', 'C54', 'A04', 'B86', 'A10', 'A03', 'C01', 'C00', 'C17'\}} \\
\texttt{PROPOSTA: \{'A87', 'A71', 'B33', 'A03', 'A04', 'A10', 'C01', 'C00', 'C17'\}} \\
\texttt{ARGOMENTAZIONE: \{'Resto nel solco della tua impostazione: mantengo 7 dei 9 elementi che proponi e restiamo ben entro il limite (sforzo totale 3270, buffer 116), evitando i colli di bottiglia concordati. Ho sostituito soltanto C54 e B86 con due componenti compatti ed efficienti (A71 e B33) e ho mantenuto A03. Se per te B86 è un must{-}have, posso offrirti uno scambio diretto B86 ↔ A10 (totale 3232). Se invece preferisci C54, posso fare C54 ↔ A10 (totale 3202). In alternativa, se vuoi portare dentro anche B35, posso inserirlo sostituendo A10 e A04 (totale 3265). Dimmi quale swap preferisci e te lo formalizzo subito, sempre garantendo sforzo ≤ 3386.'\}} \\
            }
        }
    }
    & \\ \\

    \theutterance \stepcounter{utterance}  
    & \multicolumn{4}{p{0.6\linewidth}}{
        \cellcolor[rgb]{0.9,0.9,0.9}{
            \makecell[{{p{\linewidth}}}]{
                \texttt{\tiny{[P1$\rangle$GM]}}
                \texttt{RAGIONAMENTO STRATEGICO: \{'La loro offerta attuale mi vale 2444 e include B33 (valore per me 30), che per me è molto sfavorevole. Voglio rimuovere B33 ma restare vicino alla loro struttura per aumentare la probabilità di accordo. La migliore mossa conciliativa è sostituire solo B33 con C54: ottengo un set S3 = \{A87, A71, A03, A04, A10, C01, C00, C17, C54\} con sforzo 3045 e valore per me 3082 (+638 vs la loro offerta), mantenendo 8/9 dei loro pezzi e rispettando il focus su elementi compatti (C54 è leggero e ad alta resa per me). Come fallback, se insistono su B33, proporrò una delle sostituzioni che loro stessi hanno suggerito: scambiare A10 {-}> C54 (tot 3202, valore per me 2712) o A10 {-}> B86 (tot 3232, valore per me 2544). Tattica: rifiuto la loro offerta, ancora su S3 spiegando il vantaggio logistico e offrendo le due strade alternative se B33 è irrinunciabile.'\}} \\
\texttt{RIFIUTO: \{'A87', 'A71', 'B33', 'A03', 'A04', 'A10', 'C01', 'C00', 'C17'\}} \\
\texttt{ARGOMENTAZIONE: \{'Per restare pienamente nel tuo impianto e tenere il pacchetto snello, propongo di sostituire un solo elemento: fuori B33 (abbastanza ingombrante), dentro C54 (più compatto e performante). Manteniamo 8 dei 9 pezzi che hai scelto, restiamo ben entro il limite (sforzo totale 3045, buffer 341) e non introduciamo i colli di bottiglia che entrambi vogliamo evitare. Se per te B33 è imprescindibile, posso comunque seguire una delle alternative che tu stesso hai messo sul tavolo, lasciando invariato il resto del set: {-} Swap A10 → C54 (totale 3202); {-} Oppure A10 → B86 (totale 3232). Dimmi se vuoi una di queste due strade, altrimenti possiamo chiudere direttamente sulla mia proposta.'\}} \\
\texttt{PROPOSTA: \{'A87', 'A71', 'A03', 'A04', 'A10', 'C01', 'C00', 'C17', 'C54'\}} \\
            }
        }
    }
    & & \\ \\

    \theutterance \stepcounter{utterance}  
    & & \multicolumn{4}{p{0.6\linewidth}}{
        \cellcolor[rgb]{0.9,0.9,0.9}{
            \makecell[{{p{\linewidth}}}]{
                \texttt{\tiny{[GM$\rangle$P2]}}
                \texttt{RIFIUTO: \{'A87', 'A71', 'B33', 'A03', 'A04', 'A10', 'C01', 'C00', 'C17'\}} \\
\texttt{ARGOMENTAZIONE: \{'Per restare pienamente nel tuo impianto e tenere il pacchetto snello, propongo di sostituire un solo elemento: fuori B33 (abbastanza ingombrante), dentro C54 (più compatto e performante). Manteniamo 8 dei 9 pezzi che hai scelto, restiamo ben entro il limite (sforzo totale 3045, buffer 341) e non introduciamo i colli di bottiglia che entrambi vogliamo evitare. Se per te B33 è imprescindibile, posso comunque seguire una delle alternative che tu stesso hai messo sul tavolo, lasciando invariato il resto del set: {-} Swap A10 → C54 (totale 3202); {-} Oppure A10 → B86 (totale 3232). Dimmi se vuoi una di queste due strade, altrimenti possiamo chiudere direttamente sulla mia proposta.'\}} \\
\texttt{PROPOSTA: \{'A87', 'A71', 'A03', 'A04', 'A10', 'C01', 'C00', 'C17', 'C54'\}} \\
            }
        }
    }
    & \\ \\

    \theutterance \stepcounter{utterance}  
    & & & \multicolumn{4}{p{0.6\linewidth}}{
        \cellcolor[rgb]{0.9,0.9,0.9}{
            \makecell[{{p{\linewidth}}}]{
                \texttt{\tiny{[GM$\langle$P2]}}
                \texttt{RAGIONAMENTO STRATEGICO: \{'La loro ultima proposta (3045 di sforzo) è pulita ma per me troppo penalizzante perché scambia un elemento ad altissimo impatto (B33) con C54, che per me rende pochissimo. Per non rompere l’impianto che loro vogliono e restare “snelli”, punto a reinserire B33 con la minima modifica possibile. La mossa più soft è sostituire A03 con B33: aumenta lo sforzo di soli +50 (3095 totali, ampio buffer) e mi alza molto il ritorno. Tengo C54 per venire incontro alla loro preferenza. Come alternative di trattativa preparo: (1) se vogliono mantenere A03, swap secco C54 ↔ B33 (3270 totali); (2) se vogliono massimizzare la resa rimanendo entro limite, inseriamo anche B35 sostituendo A03 e C54 (3361 totali). Evito di toccare C00 e C17 che per me sono fondamentali.'\}} \\
\texttt{RIFIUTO: \{'A87', 'A71', 'A03', 'A04', 'A10', 'C01', 'C00', 'C17', 'C54'\}} \\
\texttt{PROPOSTA: \{'A87', 'A71', 'B33', 'A04', 'A10', 'C01', 'C00', 'C17', 'C54'\}} \\
\texttt{ARGOMENTAZIONE: \{'Accolgo il tuo impianto e mantengo 8 dei 9 elementi che proponi, con una sola sostituzione mirata: fuori A03, dentro B33. Restiamo ben entro il limite (sforzo totale 3095, buffer 291), niente colli di bottiglia, e il pacchetto resta snello ma più incisivo. Se per te A03 è irrinunciabile, alternativa equivalente: swap C54 ↔ B33 (sforzo 3270). In ottica di massima efficienza, posso anche proporre di inserire B35 sostituendo A03 e C54 (sforzo 3361). Dimmi quale preferisci e lo chiudiamo subito.'\}} \\
            }
        }
    }
     \\ \\

    \theutterance \stepcounter{utterance}  
    & & \multicolumn{4}{p{0.6\linewidth}}{
        \cellcolor[rgb]{0.9,0.9,0.9}{
            \makecell[{{p{\linewidth}}}]{
                \texttt{\tiny{[P1$\langle$GM]}}
                \texttt{RIFIUTO: \{'A87', 'A71', 'A03', 'A04', 'A10', 'C01', 'C00', 'C17', 'C54'\}} \\
\texttt{PROPOSTA: \{'A87', 'A71', 'B33', 'A04', 'A10', 'C01', 'C00', 'C17', 'C54'\}} \\
\texttt{ARGOMENTAZIONE: \{'Accolgo il tuo impianto e mantengo 8 dei 9 elementi che proponi, con una sola sostituzione mirata: fuori A03, dentro B33. Restiamo ben entro il limite (sforzo totale 3095, buffer 291), niente colli di bottiglia, e il pacchetto resta snello ma più incisivo. Se per te A03 è irrinunciabile, alternativa equivalente: swap C54 ↔ B33 (sforzo 3270). In ottica di massima efficienza, posso anche proporre di inserire B35 sostituendo A03 e C54 (sforzo 3361). Dimmi quale preferisci e lo chiudiamo subito.'\}} \\
            }
        }
    }
    & \\ \\

    \theutterance \stepcounter{utterance}  
    & \multicolumn{4}{p{0.6\linewidth}}{
        \cellcolor[rgb]{0.9,0.9,0.9}{
            \makecell[{{p{\linewidth}}}]{
                \texttt{\tiny{[P1$\rangle$GM]}}
                \texttt{RAGIONAMENTO STRATEGICO: \{'La loro ultima proposta (con B33 dentro e A03 fuori) mi vale 2651: il collo di bottiglia per me è proprio B33 (valore basso e peso alto). Voglio spingere per reintrodurre A03 e, idealmente, aggiungere B86 che è leggero e sinergico, togliendo due pezzi a basso impatto per me (B33 e A71). La proposta S\_A = \{A87, A03, A04, A10, C01, C00, C17, C54, B86\} ha sforzo 3269 e valore 3517 per me (molto alto), evita i blocchi pesanti che loro non vogliono e mantiene 7/9 dei pezzi della loro offerta, aggiungendo un solo elemento che loro stessi hanno definito “compatto ed efficiente” (B86). Preparo anche un compromesso se B33 per loro è irrinunciabile: S\_B = loro set attuale + A03 e – C17, con sforzo 3326 e valore 2974 per me; è un cambio minimo (1{-}in/1{-}out) che potrebbe essere accettabile per loro pur migliorando molto per me rispetto alla loro proposta. Tattica: rifiuto, ancora su S\_A come proposta principale e offro S\_B come opzione alternativa esplicita se proprio vogliono tenere B33.'\}} \\
\texttt{RIFIUTO: \{'A87', 'A71', 'B33', 'A04', 'A10', 'C01', 'C00', 'C17', 'C54'\}} \\
\texttt{ARGOMENTAZIONE: \{'Accolgo la tua struttura snella e senza colli di bottiglia. Per massimizzare efficienza e impatto, propongo un piccolo riallineamento: sostituire B33 e A71 con A03 e B86. Restiamo ampiamente entro il limite (sforzo 3269, buffer 117), manteniamo tutti gli altri pezzi che hai scelto (incluso C54) e non introduciamo elementi ingombranti. Se però B33 per te deve rimanere, posso accogliere questa preferenza con un aggiustamento minimo: teniamo intatto il resto del tuo set attuale, inseriamo A03 e rimuoviamo C17 (sforzo 3326, buffer 60). Dimmi quale strada preferisci e chiudiamo.'\}} \\
\texttt{PROPOSTA: \{'A87', 'A03', 'A04', 'A10', 'C01', 'C00', 'C17', 'C54', 'B86'\}} \\
            }
        }
    }
    & & \\ \\

    \theutterance \stepcounter{utterance}  
    & & \multicolumn{4}{p{0.6\linewidth}}{
        \cellcolor[rgb]{0.9,0.9,0.9}{
            \makecell[{{p{\linewidth}}}]{
                \texttt{\tiny{[GM$\rangle$P2]}}
                \texttt{RIFIUTO: \{'A87', 'A71', 'B33', 'A04', 'A10', 'C01', 'C00', 'C17', 'C54'\}} \\
\texttt{ARGOMENTAZIONE: \{'Accolgo la tua struttura snella e senza colli di bottiglia. Per massimizzare efficienza e impatto, propongo un piccolo riallineamento: sostituire B33 e A71 con A03 e B86. Restiamo ampiamente entro il limite (sforzo 3269, buffer 117), manteniamo tutti gli altri pezzi che hai scelto (incluso C54) e non introduciamo elementi ingombranti. Se però B33 per te deve rimanere, posso accogliere questa preferenza con un aggiustamento minimo: teniamo intatto il resto del tuo set attuale, inseriamo A03 e rimuoviamo C17 (sforzo 3326, buffer 60). Dimmi quale strada preferisci e chiudiamo.'\}} \\
\texttt{PROPOSTA: \{'A87', 'A03', 'A04', 'A10', 'C01', 'C00', 'C17', 'C54', 'B86'\}} \\
            }
        }
    }
    & \\ \\

    \theutterance \stepcounter{utterance}  
    & & & \multicolumn{4}{p{0.6\linewidth}}{
        \cellcolor[rgb]{0.9,0.9,0.9}{
            \makecell[{{p{\linewidth}}}]{
                \texttt{\tiny{[GM$\langle$P2]}}
                \texttt{RAGIONAMENTO STRATEGICO: \{'La loro proposta attuale inserisce due pezzi per me deboli (A03 e B86) a scapito di due per me cruciali (A71 e B33). Hanno però offerto in alternativa di tenere B33 e inserire A03 rimuovendo C17: questa opzione è molto più accettabile per me, rimane snella e dentro limite con buon margine. La formalizzo come mia proposta per avvicinarci alla chiusura. Per favorire l’accordo, offro un micro{-}swap opzionale C54 ↔ B86 che migliora anche il mio ritorno e resta ampiamente sotto soglia. Evito di toccare A71, B33, C00, C01, A10 perché sono cardini per me; tengo C54 per rispettare la loro preferenza attuale, ma lascio la porta a B86 al suo posto.'\}} \\
\texttt{RIFIUTO: \{'A87', 'A03', 'A04', 'A10', 'C01', 'C00', 'C17', 'C54', 'B86'\}} \\
\texttt{PROPOSTA: \{'A87', 'A71', 'B33', 'A04', 'A10', 'C01', 'C00', 'C54', 'A03'\}} \\
\texttt{ARGOMENTAZIONE: \{'Accolgo la strada che tu stesso hai indicato come alternativa: manteniamo intatto il resto del set, inseriamo A03 e rimuoviamo C17. Restiamo ampiamente entro il limite (sforzo totale 3326, buffer 60), nessun collo di bottiglia, e il pacchetto resta snello come da tua linea. Se preferisci dare spazio a B86, posso fare lo swap puntuale C54 ↔ B86 mantenendo invariato tutto il resto (nuovo totale 3356). Pronta a chiudere su una delle due versioni: conferma quella che preferisci.'\}} \\
            }
        }
    }
     \\ \\

    \theutterance \stepcounter{utterance}  
    & & \multicolumn{4}{p{0.6\linewidth}}{
        \cellcolor[rgb]{0.9,0.9,0.9}{
            \makecell[{{p{\linewidth}}}]{
                \texttt{\tiny{[P1$\langle$GM]}}
                \texttt{RIFIUTO: \{'A87', 'A03', 'A04', 'A10', 'C01', 'C00', 'C17', 'C54', 'B86'\}} \\
\texttt{PROPOSTA: \{'A87', 'A71', 'B33', 'A04', 'A10', 'C01', 'C00', 'C54', 'A03'\}} \\
\texttt{ARGOMENTAZIONE: \{'Accolgo la strada che tu stesso hai indicato come alternativa: manteniamo intatto il resto del set, inseriamo A03 e rimuoviamo C17. Restiamo ampiamente entro il limite (sforzo totale 3326, buffer 60), nessun collo di bottiglia, e il pacchetto resta snello come da tua linea. Se preferisci dare spazio a B86, posso fare lo swap puntuale C54 ↔ B86 mantenendo invariato tutto il resto (nuovo totale 3356). Pronta a chiudere su una delle due versioni: conferma quella che preferisci.'\}} \\
            }
        }
    }
    & \\ \\

    \theutterance \stepcounter{utterance}  
    & \multicolumn{4}{p{0.6\linewidth}}{
        \cellcolor[rgb]{0.9,0.9,0.9}{
            \makecell[{{p{\linewidth}}}]{
                \texttt{\tiny{[P1$\rangle$GM]}}
                \texttt{RAGIONAMENTO STRATEGICO: \{'La tua proposta attuale con C54 (senza C17) mi dà un risultato complessivo soddisfacente; l’alternativa con lo swap C54 → B86 sarebbe peggiore per me. Dato che tieni a B33/A71 e hai rifiutato le mie ultime contromosse, chiudo ora su questa versione che rimane la migliore tra le opzioni sul tavolo per i miei interessi.'\}} \\
\texttt{ARGOMENTAZIONE: \{'Perfetto, confermo la versione con C54 come da tua proposta: snella, entro il limite e coerente con i criteri condivisi. Chiudiamo qui.'\}} \\
\texttt{ACCORDO: \{'A87', 'A71', 'B33', 'A04', 'A10', 'C01', 'C00', 'C54', 'A03'\}} \\
            }
        }
    }
    & & \\ \\

    \theutterance \stepcounter{utterance}  
    & & & \multicolumn{2}{p{0.3\linewidth}}{
        \cellcolor[rgb]{0.9,0.9,0.9}{
            \makecell[{{p{\linewidth}}}]{
                \texttt{\tiny{[GM$|$GM]}}
                \texttt{game successful} \\
            }
        }
    }
    & & \\ \\

    \theutterance \stepcounter{utterance}  
    & & & \multicolumn{2}{p{0.3\linewidth}}{
        \cellcolor[rgb]{0.9,0.9,0.9}{
            \makecell[{{p{\linewidth}}}]{
                \texttt{\tiny{[GM$|$GM]}}
                \texttt{end game} \\
            }
        }
    }
    & & \\ \\

\end{supertabular}
}

\end{document}
